\documentclass[12pt]{article}
\usepackage[T1]{fontenc}
\usepackage[utf8]{inputenc}
\usepackage{lmodern}
\usepackage{textcomp}
\usepackage{enumitem}
\usepackage{geometry}
\geometry{margin=1in}
\begin{document}
\begin{center}
Answer Key - Physics - Graduate Level Assessment 14 \
\small (Answers are ordered exactly as in the question paper.)
\end{center}

\section*{Multiple Choice}
\begin{enumerate}[label=\arabic*.]
\item \textbf{Answer:} A
\\ \textit{Options:} A) -7.89 $10^3 \mathrm{~J}$; B) -12.75 $10^3 \mathrm{~J}$; C) -8.01 $10^3 \mathrm{~J}$; D) -10.02 $10^3 \mathrm{~J}$; E) -11.32 $10^3 \mathrm{~J}$
\\ \textit{Note:} The correct answer of -7.89 x $10^3$ J is obtained by applying the Gibbs free energy change formula for isothermal processes, which accounts for the work done during the expansion of the ideal gas, leading to a negative value indicating that the process is spontaneous.
\item \textbf{Answer:} E
\\ \textit{Options:} A) Any craters that existed have been eroded through the strong winds on Io's surface.; B) Io is protected by a layer of gas that prevents any objects from reaching its surface.; C) The gravitational pull of nearby moons prevent objects from hitting Io.; D) Io's magnetic field deflects incoming asteroids, preventing impact craters.; E) Io did have impact craters but they have all been buried in lava flows.
\\ \textit{Note:} The correct answer is supported by the fact that Io experiences intense geological activity, leading to frequent volcanic eruptions that can cover and erase impact craters with new lava flows.
\item \textbf{Answer:} A
\\ \textit{Options:} A) Mass, velocity, and acceleration; B) Mass, velocity, height, and acceleration; C) Mass and velocity; D) Velocity, height, and acceleration; E) Velocity and acceleration
\\ \textit{Note:} To determine the tension in the elevator cables, one must know the mass of the elevator to calculate the gravitational force, the acceleration to account for any changes in velocity, and the velocity to understand the dynamics of the system, all of which are essential for applying Newton's second law of motion.
\item \textbf{Answer:} E
\\ \textit{Options:} A) 1000 times more; B) 50 times more; C) 5000 times more; D) 500 times more; E) 10000 times more
\\ \textit{Note:} The telescope can gather light based on the area of its aperture compared to the area of the pupil, and since the area of a circle is proportional to the square of its diameter, the ratio of the light-gathering capability is (50 cm / 0.5 cm)² = 10000, indicating the telescope gathers 10,000 times more light than the eye.
\item \textbf{Answer:} C
\\ \textit{Options:} A) reduces by 1; B) increases by 3; C) increases by 1; D) reduces by half; E) increases by 4
\\ \textit{Note:} When an element ejects a beta particle, a neutron is converted into a proton, resulting in an increase of the atomic number by 1 while the mass number remains unchanged.
\end{enumerate}

\section*{True / False}
\begin{enumerate}[label=\arabic*.]
\setcounter{enumi}{5}
\item \textbf{Answer:} True
\\ \textit{Options:} A) True; B) False
\\ \textit{Note:} Terrestrial radiation primarily consists of infrared radiation emitted by the Earth, which has a lower frequency and longer wavelength than the higher frequency and shorter wavelength ultraviolet and visible radiation emitted by the Sun.
\item \textbf{Answer:} False
\\ \textit{Options:} A) True; B) False
\\ \textit{Note:} The volcanic and tectonic history of a planet is primarily influenced by its internal heat, geological processes, and mantle dynamics, rather than the strength of its magnetic field.
\item \textbf{Answer:} True
\\ \textit{Options:} A) True; B) False
\\ \textit{Note:} The statement is true because the Fermi energy represents the maximum energy of electrons in a solid at absolute zero, and for copper, this value is indeed approximately 6.75 eV, indicating the energy level at which the highest occupied electron states exist in the material.
\item \textbf{Answer:} True
\\ \textit{Options:} A) True; B) False
\\ \textit{Note:} The power of a lens is calculated using the formula P = 1/f, where f is the focal length in meters; since the focal length for an image at 16.6 cm is 0.166 m, the power is P = 1/0.166 ≈ 6.02 diopters, indicating a miscalculation in the statement, thus making it false.
\item \textbf{Answer:} True
\\ \textit{Options:} A) True; B) False
\\ \textit{Note:} Pluto has one medium-sized satellite, Charon, and two smaller satellites, Nix and Hydra, which confirms the statement as true.
\end{enumerate}

\section*{Long Form Response}
\begin{enumerate}[label=\arabic*.]
\setcounter{enumi}{10}
\item \textbf{Answer:} To calculate the gravitational force acting on a spherical water drop with a diameter of 1.20 µm, we first determine its volume using the formula for the volume of a sphere, \( V = \frac{4}{3} \pi r^3 \), where the radius \( r = 0.6 \, \mu m = 0.6 \times 10^{-6} \, m \). This gives \( V \approx 1.44 \times 10^{-18} \, m^3 \). The mass of the drop can be found using the density of water (\( \rho \approx 1000 \, kg/m^3 \)), yielding \( m = \rho V \approx 1.44 \times 10^{-15} \, kg \). The gravitational force is then calculated using \( F = mg \), with \( g \approx 9.81 \, m/s^2 \), resulting in a force of approximately \( 10.00 \times 10^{-15} \, N \). This small magnitude of gravitational force highlights the challenges faced by such tiny particles in balancing gravitational and electric forces, which is crucial in understanding phenomena such as electrostatic suspension and the behavior of aerosols.
\\ \textit{Note:} correct because it accurately employs the formula for the volume of a sphere to determine the water drop's volume, calculates its mass using water's density, and thus successfully derives the gravitational force acting on the drop, demonstrating an understanding of the relationship between volume, mass, and gravitational force in the context of a microscopic object.
\item \textbf{Answer:} To analyze the effect of a 100 dyne force applied for 10 seconds on a free particle with a mass of 20 grams, we first calculate the acceleration using Newton's second law, \( F = ma \). The mass in grams must be converted to kilograms for consistency in SI units, which gives us \( m = 0.02 \) kg. The force of 100 dyne is equivalent to \( 0.01 \) N (since \( 1 \, \text{dyne} = 10^{-5} \, \text{N} \)). Thus, the acceleration \( a \) is \( a = F/m = 0.01 \, \text{N} / 0.02 \, \text{kg} = 0.5 \, \text{m/s}^2 \). Next, we determine the final velocity after 10 seconds using the equation \( v = u + at \), where \( u = 0 \) (initial velocity). This gives \( v = 0 + (0.5 \, \text{m/s}^2)(10 \, \text{s}) = 5 \, \text{m/s} \). The kinetic energy (KE) is calculated using the formula \( KE = \frac{1}{2} mv^2 \). Substituting the values, we find \( KE = \frac{1}{2} (0.02 \, \text{kg})(5 \, \text{m/s})^2 = \frac{1}{2} (0.02)(25) = 0.25 \, \text{J} \), which converts
\\ \textit{Note:} correct because it accurately applies Newton's second law to determine the acceleration of the particle from the applied force, converts units appropriately for consistency, and sets the stage for further calculations of kinetic energy based on the resulting motion.
\end{enumerate}

\vfill
\noindent\textit{End of Answer Key}
\end{document}